% !PRINTED 454
\weberpage{454}

%<<CONT
double point, point of inflection, cusp, point of contact of a double tangent, etc.

Among the forms $\varphi$, those are now above all of importance to us which remain unchanged under the whole group, the invariants, of which there are here, since the group is simple, only absolute ones (\S.\,40). If $\Phi\,(x_1, x_2, x_3)$ is such a form, then the curve represented by the equation $\Phi = 0$ shall be called an invariant curve of the group. There are 168 mappings of the plane onto itself under which all points of the curve pass into points of the same curve. If any point lies on this curve, then all its image points also lie on it.

In general the number of points of the curve thus connected with one another is 168, in any case not greater. If it is smaller, then the points must be \gesperrt{poles}, and the number of image points is a divisor of 168 (namely 24 for the $P_7$, 21 for the $P_8$, 28 for the $P_6$, 42 for the $P_4$, 56 for the $P_3$, and 84 for a system of $P_2$ belonging together). If a pole of one of these kinds lies on the curve, then all poles of the same kind lie on it.

\begin{center}
\S.\,122.\\
\gesperrt{The first invariant of the group $G_{168}$ and the fundamental curve.}
\end{center}

In order now to form the invariants of our group, we seek out forms of the three variables $x_1, x_2, x_3$ that remain unchanged by the application of the three substitutions $\chi, \tau, \omega$ (\S.\,115). This suffices, since the whole group can be composed from these three substitutions (\S.\,72). Now $\chi$ is a cyclic permutation of the three variables $x_1, x_2, x_3$, and $\tau$ denotes the substitution

\[
\begin{pmatrix} x_1, & x_2, & x_3 \\ \varepsilon x_1, & \varepsilon^2 x_2, & \varepsilon^4 x_3 \end{pmatrix}.
\]

If, then, a term $x_1^{h_1} x_2^{h_2} x_3^{h_3}$ occurs in an invariant, in which $h_1, h_2, h_3$ are non-negative integers, then invariability under $\tau$ requires that

\weq{1}{h_1 + 2h_2 + 4h_3 \equiv 0\ (\mathrm{mod}\ 7)}
%CONT>>
